% Responsable: Martín (investigación y contexto)
% Tesis: el Estado colombiano focaliza con información imperfecta; la pregunta
% del paper es cuánto cuesta esa imperfección y cómo se reparte el gasto.
\section{Introducción}\label{sec:intro}

\pendiente{Contexto: IPM oficial en Colombia y Antioquia (cifras DANE 2025).}

\pendiente{Problema: el gasto social se asigna sin saber exactamente quién es
pobre (errores de inclusión y exclusión).}

\pendiente{Aporte: adaptamos los cuatro modelos de \citet{hamie2024} a Antioquia
y medimos el costo de pasar de información perfecta (Modelos 1--2) a
focalización geográfica (Modelo 3) y mixta (Modelo 4).}

\pendiente{Pregunta de investigación y estructura del documento
(Sección~\ref{sec:literatura} a Sección~\ref{sec:conclusiones}).}
